\documentclass[11pt,reqno]{amsart}
\usepackage{amsmath,amssymb,amsthm}
\usepackage[margin=1.1in]{geometry}
\usepackage{graphicx}
\usepackage{array,longtable}
\usepackage{hyperref}
\hypersetup{colorlinks=true,linkcolor=blue,citecolor=blue,urlcolor=blue}

\theoremstyle{plain}
\newtheorem{theorem}{Theorem}[section]
\newtheorem{lemma}[theorem]{Lemma}
\newtheorem{corollary}[theorem]{Corollary}
\newtheorem{proposition}[theorem]{Proposition}
\theoremstyle{definition}
\newtheorem{definition}[theorem]{Definition}
\newtheorem{remark}[theorem]{Remark}
\newtheorem{question}[theorem]{Question}
\newtheorem{computation}[theorem]{Computation}
\newtheorem{conjecture}[theorem]{Conjecture}

\newcommand{\R}{\mathbb{R}}
\newcommand{\Z}{\mathbb{Z}}
\newcommand{\Q}{\mathbb{Q}}
\newcommand{\N}{\mathbb{N}}
\DeclareMathOperator{\Area}{Area}
\newcommand{\sF}{\textup{\textsf{F}}}
\newcommand{\sW}{\textup{\textsf{W}}}
\newcommand{\sS}{\textup{\textsf{S}}}
\DeclareUrlCommand{\lean}{\urlstyle{tt}}
\newcommand{\rec}[1]{\textup{\texttt{#1}}}
\emergencystretch=2em\sloppy

\title[Release 6.0]{Erd\H{o}s Problem 634, release 6.0: the W-gap theorem, the notch and band families, tile counts by explicit tilings, and the census up to $10^6$}
\author{Vico Bonfioli}
\date{18 November 2026 (release 6.0; not refereed)}
\newcommand{\TODOCOUNT}[1]{\textbf{[\,TODO: final count\,]}}

\begin{document}
\begin{abstract}
Erd\H{o}s Problem~634 asks for the integers $N$ for which some triangle can be cut into $N$ congruent triangles
(\emph{tile counts}). \emph{The problem is not solved here.} Since release v5.0 we add: a Lean theorem that a family of
triangles $W_M(e,f)$ (defined in Section~\ref{sub:intro:notation}) is not tiled by the tile $\tau(e,f)$ for coprime
$2\le e<f$ when $f\le Me$ and $M(f-e)<f$, with per-value Lean instances for the $5{,}593$
values it excludes (all $5{,}593$ checked by the Lean kernel); an explicit threshold and a band theorem (Lean) for isosceles targets
with base angles $\alpha+\beta$; three structural theorems (Lean); a refutation of a conjectured foot property;
and $132$ further tile counts by explicit tilings (exact programs, not Lean). The undecided values $N\le10^6$ fall
from $62{,}372$ to $52{,}045$. The bottom rung of $W$ and the thresholds for $T_M$ stay open. All checks are internal; no outside reader has checked anything.
\end{abstract}
\maketitle

\section{Overview: the problem, what is new, and how to read this paper}\label{sec:intro}

\emph{Roadmap.} This section states the problem in plain words, says what is new and what is not claimed, explains the
labels, fixes the notation (Table~\ref{tab:notation}), gives the status of the problem (Table~\ref{tab:status}) and
lists the results (Table~\ref{tab:glance}). Sections~\ref{sec:wgap}--\ref{sec:band} contain the theorems,
Section~\ref{sec:tilings} the explicit tilings, Section~\ref{sec:census} the census; the credits and the disclosure are at the end.

\subsection{The problem}\label{sub:intro:problem}
A positive integer $N$ is a \emph{tile count} if some triangle can be cut into $N$ pairwise congruent triangles
(the tiles may be mirror images of each other; the cut need not be edge-to-edge). Erd\H{o}s Problem~634 asks which $N$ are
tile counts. The hard direction is to \emph{prove that
a given $N$ is not} a tile count. We say that such an $N$ is \emph{excluded}. An $N$ is \emph{decided} if it is
known to be a tile count or is excluded, and \emph{undecided} otherwise.

\subsection{What is new, and what is not claimed}\label{sub:intro:new}
\textbf{New since release v5.0 (1 November 2026).}
(1)~A theorem, checked in Lean, that a certain family of large triangles $W_M(e,f)$ cannot be tiled by the tile
$\tau(e,f)$ in a range of parameters (the \emph{W-gap theorem}); it excludes $5{,}593$ values of $N$.
(2)~A construction and a theorem, checked in Lean, for isosceles targets (the \emph{notch family} and the \emph{band theorem}),
which give $1{,}410$ new tile counts. (3)~Three structural statements about every tiling of $W_M$ and $T_M$.
(4)~A refutation of a conjectured property (the \emph{foot property}). (5)~$132$ further tile counts by explicit tilings,
verified by exact programs, not by Lean. (6)~The census: of the non-squarefree $N\le10^6$, the number of undecided
values fell from $62{,}372$ (v5.0) to $52{,}045$.

\textbf{Not claimed.} The problem is not solved: Table~\ref{tab:status} shows what is still open. We do not claim that any
undecided $N$ is or is not a tile count. The bottom rung of $W$, the thresholds for $T_M$ and $\theta_t$, and the
family $\gamma=2\alpha$ stay open. All checking is internal to the project; no reader outside the project has checked any
statement. Nothing here bears on the prime case.

\subsection{How to read the labels}\label{sub:intro:labels}
Every result carries one mark.
\begin{itemize}\setlength{\itemsep}{0pt}
\item \sF{} \textbf{(Lean)}: a theorem of the Lean~4 library, checked by the Lean kernel, with no \texttt{sorry} and no axiom
beyond \texttt{propext}, \texttt{Classical.choice}, \texttt{Quot.sound}, and compared with the printed statement by an
auditor inside the project.
\item \sF$^\ast$: as \sF{}, but compared with the printed statement by one reader only.
\item \sW{} \textbf{(hand)}: a written proof, not formalized.
\item \sS{} \textbf{(program)}: the output of a computer search or count, with exact verification but without a Lean
certificate.
\end{itemize}
Nothing in this paper is marked as checked by a person outside the project. ``Lean-checked'' and ``kernel-checked'' both mean
\sF{} or \sF$^\ast$; a tiling that is only \sS{} is called \emph{machine-verified}. Lean names and file positions are
printed after each theorem so that a reader can look them up in the formalization repository.

\subsection{Notation}\label{sub:intro:notation}
Throughout, $e<f$ are coprime positive integers and $\tau=\tau(e,f)$ is the \emph{tile}, a triangle with sides
$a=ef$, $b=f^2-e^2$, $c=f^2$. Its angles are $\alpha,\beta,\gamma$, opposite $a,b,c$. They satisfy
$\gamma=2\alpha+\beta$ and $3\alpha+2\beta=\pi$. In words: $c$ is the longest side and $\gamma$ the largest angle, and the three angles are in the proportions fixed by
$\gamma=2\alpha+\beta$. (Section~\ref{sec:tilings} uses a different
family of tiles with an angle $2\pi/3$, and re-uses the letters $a,b,c$ there; it says so.)

\begin{table}[h]
\centering\small
\begin{tabular}{@{}p{3.1cm}p{9.6cm}@{}}
\hline
Symbol & Meaning\\\hline
$\tau(e,f)$ & the tile, sides $a=ef$, $b=f^2-e^2$, $c=f^2$, angles $\alpha,\beta,\gamma$ opposite $a,b,c$\\
$N$ & number of tiles; $N$ is a \emph{tile count} if some triangle is cut into $N$ congruent triangles\\
$M$ & integer scale factor ($M\ge1$): sides of the targets below are multiplied by $M$\\
$W_M(e,f)$ & the triangle with angles $\beta,\,2\alpha,\,\alpha+\beta$ and sides $Mf(f^2-e^2)$, $Me(2f^2-e^2)$, $Mf^3$ (opposite those angles); it needs $(2f^2-e^2)M^2$ tiles\\
$T_M(e,f)$ & the isosceles triangle with apex angle $3\alpha$ (so base angles $\beta$), legs $Mf^3$, base $Me(3f^2-e^2)$\\
$\theta_t(e,f)$ & the isosceles triangle with apex $\alpha$, base angles $\alpha+\beta$, legs $tf(f^2-e^2)$, base $te(f^2-e^2)$; it needs $(f^2-e^2)t^2$ tiles\\
$\Theta(e,f)$ & the set of $t\ge1$ for which $\theta_t(e,f)$ is tiled by $\tau(e,f)$\\
$\alpha_g(e,f)$ & the isosceles triangle with base angles $\alpha$ (Section~\ref{sub:band:setting})\\
$f-e$, $M(f-e)$ & the gap of the tile; the \emph{W-window} is $f\le Me$ and $M(f-e)<f$\\
\sF,\ \sF$^\ast$,\ \sW,\ \sS & the marks of Section~\ref{sub:intro:labels}\\\hline
\end{tabular}
\caption{Notation. A triangle is \emph{tiled} (by $\tau$) if it is the union of congruent copies of $\tau$ with disjoint interiors.}
\label{tab:notation}
\end{table}

\subsection{Status of the problem}\label{sub:intro:status}
The census (Section~\ref{sec:census}) concerns the non-squarefree integers $N\le10^6$ ($392{,}074$ values); squarefree $N$
are treated by other results of the project and are not counted here.
\begin{table}[h]
\centering\small
\begin{tabular}{@{}p{7.4cm}rp{3.6cm}@{}}
\hline
Statement & Count & Status\\\hline
Non-squarefree $N\le10^6$ (the census universe) & $392{,}074$ & ---\\
Decided (tile count or excluded) & $340{,}029$ & decided\\
Undecided $N\le10^6$ & $52{,}045$ & open\\
Undecided $N\le10^6$ at v5.0 (1 Nov.\ 2026) & $62{,}372$ & fell by $10{,}327$\\
Undecided $N\le10^5$ / $\le10^4$ / $\le3{,}000$ & $5{,}115$ / $462$ / $111$ & open\\
Undecided $N\le1000$ & $27$ & open\\
Every $N\le200$ & --- & decided\\
Values decided by the $W$-theorems (this release) & $5{,}593$ & decided (Lean)\\
Values newly shown to be tile counts by explicit tilings & $132$ & decided (program)\\
Conjectures T, theta, gamma (Section~\ref{sub:cen:conj}) & $3$ & conjecture, open\\\hline
\end{tabular}
\caption{Status of the problem as measured by the census of 2026-11-16g.}
\label{tab:status}
\end{table}

\subsection{Results at a glance}\label{sub:intro:glance}
Open: the bottom rung of $W$, the thresholds for $T_M$, $\theta_t$, and the $\gamma=2\alpha$ family (Section~\ref{sub:cen:conj}).
\begin{table}[h]\centering\footnotesize
\begin{tabular}{@{}p{4.9cm}p{0.9cm}p{2.6cm}p{4.5cm}@{}}
\hline
Result & Label & Where & Lean name\\\hline
W-gap theorem: $W_M(e,f)$ not tiled for $f\le Me$, $M(f-e)<f$ & \sF & Thm.~\ref{thm:wgap} & \texttt{no\_tiling\_W\_gap}\\
Rungs 2--5 (earlier, weaker) & \sF & Cor.~\ref{cor:wrung} & \texttt{no\_tiling\_W\_rung$k$}\\
Per-value instances $\neg\,\mathrm{TileCount}\,N$ for the $5{,}593$ $W$-values & \sF & \S\ref{sub:wgap:light} & per file; $5{,}593$ of $5{,}593$ done\\
Long run at the $2\alpha$ corner of $W_M$ & \sF$^\ast$ & Thm.~\ref{thm:wrun} & \texttt{w\_corner\_run\_long}\\
Bottom-rung corner words & \sF$^\ast$ & Thm.~\ref{thm:wbottom} & \texttt{w\_bottom\_rung\_words}\\
One apex leg of $T_M$ has a long run & \sF$^\ast$ & Thm.~\ref{thm:tapex} & \texttt{t\_apex\_one\_leg\_long}\\
Notch seed family, $t\ge t_0(e,f)$ & \sF & Thm.~\ref{thm:notch} & \texttt{thetaTiled\_all}\\
Case $e=1$ exact: $t\ge2f$ & \sF & Thm.~\ref{thm:e1} & \texttt{thetaTiled\_e1\_iff}\\
Alpha hub & \sF & Prop.~\ref{prop:hub} & \texttt{alphaTiled\_of\_thetaTiled}\\
Band theorem; odd-$e$ corollary; $1{,}410$ instances & \sF & Thm.~\ref{thm:band}, Cor.~\ref{cor:odde} & \texttt{thetaTiled\_band}\\
FOOT-U is false ($T_3(1,2)$, $T_4(1,2)$) & \sS & Prop.~\ref{prop:foot} & ---\\
$132$ tile counts by explicit tilings & \sS & \S\ref{sec:tilings}; companion data paper & ---\\
Census: undecided $N\le10^6$: $62{,}372\to52{,}045$ & \sS & \S\ref{sec:census} & ---\\
Hand proof of the W-gap theorem (not used by Lean) & \sW & \S\ref{sub:wgap:thm} & ---\\\hline
\end{tabular}
\caption{Results of this release, with marks and Lean names.}\label{tab:glance}
\end{table}

\section{The W-gap theorem}\label{sec:wgap}

\emph{Roadmap.} We define the triangles $W_M(e,f)$ and link them to tile counts (\ref{sub:wgap:setting}); state the
W-gap theorem and sketch its proof (\ref{sub:wgap:thm}); state weaker rung versions and three structural theorems
(\ref{sub:wgap:struct}); describe the per-value Lean instances (\ref{sub:wgap:light}); report a refuted conjecture,
the foot property (\ref{sub:wgap:foot}); and list what the theorem does not cover (\ref{sub:wgap:limits}).

\begin{remark}[Standing of this section]\label{rem:wgap:status}
The theorems of this section are theorems of the Lean library (mark \sF), with the standard axioms
\texttt{propext}, \texttt{Classical.choice}, \texttt{Quot.sound}, and no \texttt{sorry}. They have been audited
inside the project only: an auditor who was not the builder compared each formal statement with the statement
printed here, hypothesis by hypothesis (work items \rec{wgap\_2026-11-16}, \rec{wrung\_2026-11-16},
\rec{wgap\_mid\_2026-11-16}, \rec{wgap\_rlgen\_2026-11-16}). \emph{No reader outside the project has checked
any of them.} The three structural theorems of Section~\ref{sub:wgap:struct} were built after these audits and
had one statement-match audit on 16 November 2026 (one reader, inside the project; work item
\rec{NEVERTRIED\_T\_2026-11-16/V20.md}), which found the formal statements to match the printed ones. The section excludes certain values of $N$ as tile counts; it contains no
construction, and it does not decide the problem.
\end{remark}

\subsection{Setting}\label{sub:wgap:setting}
Let $2\le e<f$ be coprime integers. Put
\[
 a=ef,\qquad b=f^2-e^2,\qquad c=f^2,
\]
and let $\tau=\tau(e,f)$ be the triangle with sides $a,b,c$ (opposite its angles $\alpha,\beta,\gamma$); then
$\gamma=2\alpha+\beta$ and $3\alpha+2\beta=\pi$. For an integer $M\ge1$ let $W_M(e,f)$ be the triangle with sides
\[
 M\,f\,(f^2-e^2),\qquad M\,e\,(2f^2-e^2),\qquad M\,f^3,
\]
whose angles are $\beta$, $2\alpha$, $\alpha+\beta$ (in this order, opposite the three sides). Its area is
$(2f^2-e^2)M^2$ times the area of $\tau$. We say that $W_M(e,f)$ is \emph{tiled} if it is the union of
$(2f^2-e^2)M^2$ triangles congruent to $\tau$ with pairwise disjoint interiors; the tiles may be mirror images,
and the tiling is not assumed edge-to-edge. In Lean this is
\texttt{Erdos634.Draft.CollarG.WTilable e f M} (\lean{lean/Erdos634/Draft/CollarG.lean:615}): there is a
\texttt{CongruentDissection} of $(2f^2-e^2)M^2$ tiles whose target has the three displayed side lengths and whose
model has sides $(ef,\,f^2-e^2,\,f^2)$, where \texttt{HasSides} allows any vertex labelling.

\emph{Link to tile counts} (Lean: \rec{s4\_wTilable}, \lean{lean/Erdos634/Draft/WGapLink1.lean:36}; \sF, audited
only as part of the census chain of Section~\ref{sec:census}). Suppose a congruent dissection of $N$ triangles has a
target with angles $(2\alpha,\beta,\alpha+\beta)$, for a tile with $3\alpha+2\beta=\pi$, $\alpha/\pi$ irrational,
and sides $a=rc$, $b=(1-r^2)c$ with $r$ rational. Then there are coprime $0<e<f$ and $M\ge1$ with
$N=(2f^2-e^2)M^2$ and $W_M(e,f)$ tiled. In words: a tile count of this kind is always of the form
$(2f^2-e^2)M^2$, so the W-gap theorem removes the representations of $N$ of this type that lie in its window (a value is excluded only
when all its representations are, Section~\ref{sub:wgap:limits}).

\subsection{The theorem}\label{sub:wgap:thm}

\begin{theorem}[W-gap theorem; \sF]\label{thm:wgap}
Let $e,f,M$ be integers with $\gcd(e,f)=1$, $2\le e<f$, $f\le Me$ and $M(f-e)<f$. Then $W_M(e,f)$ is not tiled
by copies of $\tau(e,f)$.
\end{theorem}

\noindent\emph{In words.} In the stated window, one cannot fill $W_M(e,f)$ with copies of the tile, however they are
arranged. \emph{Smallest example.} Take $(e,f,M)=(2,3,2)$: the tile has sides $(6,5,9)$, the window conditions read $3\le4$
and $2<3$, and $W_2(2,3)$ has sides $30$, $56$, $54$; the theorem says it is not cut into $(2\cdot 9-4)\cdot 4=56$ copies of the tile.
For contrast, $W_3(2,3)$ (with $126$ tiles) has $M(f-e)=3=f$, lies just outside the window, and does have a
kernel-checked tiling.

\noindent Lean: \texttt{Erdos634.Draft.WGapRung.no\_tiling\_W\_gap}
(\lean{lean/Erdos634/Draft/WGapStation4.lean:70}),
\begin{center}
\begin{minipage}{0.95\textwidth}\small\ttfamily no\_tiling\_W\_gap (e f M : N) (hco : Nat.Coprime e f)\\ (he : 2 <= e) (hef : e < f) (hfM : f <= M*e)\\ (hMf : M*(f-e) < f) : not WTilable e f M\end{minipage}
\end{center}
There is no further hypothesis (the theorem is unconditional; an earlier version of the project had it
conditional on a statement \rec{WGapMainStmt}, now discharged). The proof is the module chain
\texttt{WGapStation1..4}, \texttt{WGapMainRed}, \texttt{WGapFinal}, \texttt{WGapRung7..9}, with 730 project
modules in its import closure; no module contains \texttt{sorry}, \texttt{axiom}, \texttt{native\_decide},
\texttt{unsafe} or \texttt{implemented\_by}. The hypotheses are satisfiable together (for example
$(e,f,M)=(2,3,2)$, $(5,7,3)$, $(100,101,2)$); the geometric preliminaries of the proof are satisfied by the real
tiling of $W_3(2,3)$ (a kernel-checked tiling, outside the window).

\paragraph{Proof (sketch; \sW; the Lean proof does not use this text).}
The hand argument has two halves, a structure theorem and a count, which contradict each other.
\paragraph{Step 1 (structure theorem, for every $M\ge1$).}
The lengths of the sides enter only through the hypothesis $|s'|>(f-1)c$ (see the errata below). Let $AQ$ be the side of length $Mf(f^2-e^2)$ that is
adjacent to the corner of angle $2\alpha$. In a lattice chart at that corner, in which the tile edges $c$ and $b$
span the first lattice, the tiles along $AQ$ begin with a run of $f$ consecutive tiles with their $c$-edge on
$AQ$ (alternative~(i) of the corner-run theorem, a statement of the boundary-run method of Beeson--Ro\c{s}u),
because alternative~(ii) is excluded by a joint induction on the rows and anti-diagonals of the lattice (the
lemmas called AD, U3, the wall lemma, stations, the recursion RL and the climb in PROOF.md \S\S3--10). Hence the
number $z$ of $c$-edges on $AQ$ satisfies $z\ge f$.
\paragraph{Step 2 (count).} Since
$x\,ef+y\,(f^2-e^2)+z\,f^2=Mf(f^2-e^2)$, the window $f\le Me$, $M(f-e)<f$ forces $z<f$ (a word-count argument in
the semigroup $\langle a,b,c\rangle$, \rec{SplusThick.C\_AQ'}).
\paragraph{Step 3 (conclusion).} Steps 1 and 2 give $z\ge f$ and $z<f$; hence no tiling exists.
%%ERRATA-BEGIN
\paragraph{Sources and errata of the hand proof.}
The written proof is \rec{WGAP\_PROOF\_2026-11-15/PROOF.md}, with an errata block from the formalization
(findings F5, F7--F11). F5 and F7--F10 are bookkeeping of the invariant and change no statement. F11 \emph{does}
correct a stated clause: in lemma U3 (and its variant U3-ST) the far-side tile beside the last lattice $a$-edge
need not show the angle $\beta$ at the end point of the run (clause (U3-c) keeps only that the tile beside each
$a$-step of the chord has exactly that $a$-edge), and the companion statement (D) holds only for $j\le J-1$
(resp.\ $j\le H-1$); the places that used the dropped part were re-read and derive it locally, so the corrected
induction still closes, but the earlier hand draft was wrong there. A further finding of the hand proof (U3-1) adds
the hypothesis $|s'|>(f-1)c$ to the lemma on the corner at the end of the run; it holds for $W_M$ (the side has
length $Mf^3$) and is open for general triangles. The formal theorem does not use any of these statements.
The hand proof has been read only by its author and a reader of the same model family; it is not claimed as an
independently checked proof, and the Lean theorem does not depend on it.

\subsection{Corollaries and structural theorems}\label{sub:wgap:struct}

\begin{corollary}[Rungs; \sF]\label{cor:wrung}
Under the hypotheses of Theorem~\ref{thm:wgap}, if moreover $M(f-e)\le e+j$ for one $j\in\{2,3,4,5\}$, then
$W_M(e,f)$ is not tiled. (These are weaker than Theorem~\ref{thm:wgap}; they were proved first by a shorter
induction, \texttt{no\_tiling\_W\_gap\_of\_sigma\_le}, and the lists of $N$ they decide are in Section~\ref{sec:census}.)
\end{corollary}
\noindent Lean: \texttt{no\_tiling\_W\_rung2} (\lean{WGapRung1.lean:110}), \texttt{no\_tiling\_W\_rung3}
(\lean{WGapRung3.lean:41}), \texttt{no\_tiling\_W\_rung4} (\lean{WGapRung5.lean:40}),
\texttt{no\_tiling\_W\_rung5} (\lean{WGapRung7.lean:56}); same arguments as Theorem~\ref{thm:wgap}, one more
hypothesis \texttt{hk : M*(f-e) <= e + j}.

Two statements concern every tiling and decide no $N$ by themselves. They are the geometric half of the proof,
extracted. \emph{Their statements had one statement-match audit (16 November 2026; one reader, inside the project).}

\begin{theorem}[Long run at the $2\alpha$ corner, any $M$; \sF, statement-match audited once, internally]\label{thm:wrun}
Let $e<f$ be coprime, $1\le e$, $M\ge1$, and let $D$ be a dissection of $W_M(e,f)$ into copies of $\tau(e,f)$
with sides as above (any $N$). Then there are integers $x,y,z\ge0$ such that the side adjacent to the corner of
angle $2\alpha$ carries the long-run word $\texttt{WLongRunWord}\ D\ M\ e\ f\ x\ y\ z$, with
$x\,ef+y\,(f^2-e^2)+z\,f^2=Mf(f^2-e^2)$ and $z\ge f$.
\end{theorem}
\noindent\emph{In words.} Whatever the tiling, the side next to the $2\alpha$ corner starts with a run of at least $f$ tiles
whose longest side $c$ lies on it, and the whole side is a sum of $a$'s, $b$'s and $c$'s.\\
\noindent Lean: \texttt{w\_corner\_run\_long} (\lean{WCornerStruct1.lean:72}). Neither $f\le Me$ nor $2\le e$ is
used. Control: \texttt{no\_tiling\_W\_gap'} (\lean{WCornerStruct1.lean:138}) re-derives Theorem~\ref{thm:wgap}
from this theorem and the count of $c$-tiles.

\begin{theorem}[The bottom rung; \sF, statement-match audited once, internally]\label{thm:wbottom}
Let $e<f$ be coprime, $2\le e$, $f\le Me$ and $(M-1)(f-e)<f\le M(f-e)$. In every dissection of $W_M(e,f)$ into
copies of $\tau(e,f)$ the long-run word of Theorem~\ref{thm:wrun} has $y=0$ and either $x=z=M(f-e)$, or
$x=M(f-e)-f$ and $z=M(f-e)+e$.
\end{theorem}
\noindent\emph{In words.} On the bottom rung the side next to the $2\alpha$ corner uses no $b$-edge and has only two possible
counts of $a$- and $c$-edges; this is not yet a contradiction.\\
\noindent Lean: \texttt{w\_bottom\_rung\_words} (\lean{WCornerStruct2.lean:145}). This is a statement about counts; it
does not exclude any tiling.

\begin{theorem}[One leg of the apex of $T_M$ has a long run; \sF, statement-match audited once, internally]\label{thm:tapex}
Let $1\le e<f$ be coprime, $M\ge1$, and let $T_M(e,f)$ be the isosceles triangle with legs $Mf^3$ and base
$Me(3f^2-e^2)$ (apex angle $3\alpha$) tiled by copies of $\tau(e,f)$. Then at least one of the two legs from the
apex satisfies \texttt{TLegLong}, that is, alternative~(i) of the corner-run theorem: it carries $f$ consecutive
tiles of the type $C_1$, or $f$ consecutive tiles of the type $C_2$. The statement does not say that alternative~(ii) is absent on the other leg.
\end{theorem}
\noindent\emph{In words.} At the apex of $T_M$, at least one of the two legs begins with a run of $f$ tiles of the same type.\\
\noindent Lean: \texttt{t\_apex\_one\_leg\_long} (\lean{TApexStruct1.lean:164}). It is the scale-$M$ form of
Beeson--Ro\c{s}u's Theorem~37 for the apex.

\subsection{Per-value instances of the exclusion}\label{sub:wgap:light}
The census (Section~\ref{sec:census}) uses the link of Section~\ref{sub:wgap:setting} one value at a time: for each
$N$ whose open representations are all of the type $W_M(e,f)$ in the window of Theorem~\ref{thm:wgap} (or its rungs), a
Lean module proves $\neg\,\texttt{TileCount}\ N$ by applying \texttt{not\_tileCount\_of\_CNzFixW} to the decided
representation table of $N$. These instances cover the $5{,}593$ values that the $W$ exclusions remove from the census
($77+65+98$ by rungs $2$--$4$, $5{,}351$ by Theorem~\ref{thm:wgap}, $2$ by \texttt{CNzFixW}). They are kernel-checked per
value in Lean (generator, manifest and axioms replay in the repository): all $5{,}593$ of $5{,}593$ instances have been
built, and each depends only on the three standard axioms (\texttt{propext}, \texttt{Classical.choice},
\texttt{Quot.sound}); a comment-aware scan of the import closure found no \texttt{sorry}, \texttt{native\_decide},
\texttt{ofReduceBool} or \texttt{axiom}. The $163$ values of rungs $3$ and $4$ and $74$ of the $77$ values of rung $2$
pass the same predicate as the values of the main theorem, so the Lean predicate agrees with its independent
transcription on every one of them.

\subsection{A refuted conjecture: the foot property}\label{sub:wgap:foot}
In the search for a route to the thin part of the $T_M$ wall, the following was formulated. Put $n=Me(2f^2-e^2)$; on the
base of $T_M(e,f)$ there is a corner $B$ and a tile vertex $Q$ with $|QB|=n$. The \emph{foot property} FOOT$^\ast$ says that
the first tile at $Q$ along $[Q,B]$ is not the tile type with the angle $\gamma$ at $Q$, together with a companion
condition on the second tile. FOOT$^\ast$ together with $Me<f$ gives a contradiction by a hand count (not formalized), so the
conjecture FOOT-U (every tiling of every $T_M(e,f)$ has FOOT$^\ast$) would have excluded the thin part.

\begin{proposition}[FOOT-U is false; \sS]\label{prop:foot}
$T_3(1,2)$ ($N=99$) and $T_4(1,2)$ ($N=176$) have tilings without FOOT$^\ast$ at either base corner.
\end{proposition}
\noindent\emph{In words.} The hoped-for property (``the first tile at $Q$ is not of type $\gamma$'') fails in actual
tilings, the smallest cases being $N=99$ and $N=176$.\\
\noindent The tilings were produced by the project's engine, and are checked as valid dissections by an exact rational
checker (not by Lean). In a search of $400{,}000$ tilings of $T_3(1,2)$, $56{,}972$ violate FOOT$^\ast$; of $260{,}000$ tilings of
$T_4(1,2)$, $254{,}800$ do. All $252$ tilings of $T_2(1,2)$ and all sampled instances with $M\le f$ satisfy it. This
closes the route through FOOT$^\ast$ for $M>f$; it is not evidence about whether any open $T_M(e,f)$ is tiled, and the $T$-threshold
statement (Conjecture~\ref{conj:T}) is unaffected.

\subsection{What is not covered}\label{sub:wgap:limits}
\begin{enumerate}
\item The \emph{bottom rung}, $M(f-e)\ge f>(M-1)(f-e)$ with $f\le Me$, is outside Theorem~\ref{thm:wgap}. Theorem~\ref{thm:wbottom}
 gives its two corner words only. Whether $W_M(e,f)$ is tiled there is open; the bottom rung carries 2,529 open
 representations on the list of Section~\ref{sec:census}.
\item The cases $e=1$ and $Me<f$ (the latter excluded earlier by the side theorem of the base-$\beta$ program),
 and the range of large $M$ where $W_M(e,f)$ is known to be tiled, are outside the window; the theorem says
 nothing about them.
\item Targets other than $W_M(e,f)$, in particular $T_M(e,f)$ (apex $3\alpha$) and the isosceles targets of
 Section~\ref{sec:band}: Theorem~\ref{thm:tapex} is all that is proved for $T_M$.
\item The decision of a value of $N$ requires every representation of $N$ as a tile count to be excluded; the
 theorems exclude $W$-representations only (Section~\ref{sec:census}).
\end{enumerate}

\section{Isosceles targets with base angles $\alpha+\beta$: the notch family and the band theorem}\label{sec:band}

\emph{Roadmap.} This section constructs tilings (it excludes nothing). We define the isosceles targets $\theta_t(e,f)$,
state the notch seed family with its explicit threshold $t_0(e,f)$ (\ref{sub:band:notch}), the exact case $e=1$,
the band theorem and its odd-$e$ corollary (\ref{sub:band:band}), the $1{,}410$ instances (\ref{sub:band:inst}), and what is not claimed.

\begin{remark}[Standing]\label{rem:band:status}
All theorems here are theorems of the Lean library (mark \sF; standard axioms; no \texttt{sorry}), audited
inside the project only (work items \rec{thetanotch\_2026-11-14}, \rec{thetanotch4\_2026-11-14},
\rec{thetaband\_2026-11-15}, \rec{thetabandinst\_2026-11-15}); no reader outside the project has compared the formal
statements with the printed ones. They construct tilings: they add tile counts and exclude nothing.
\end{remark}

\subsection{Setting}\label{sub:band:setting}
Let $1\le e<f$ and put $b=f^2-e^2$, $\tau=\tau(e,f)=(ef,\,b,\,f^2)$ as in Section~\ref{sub:wgap:setting}. For an integer $t\ge1$ let
$\theta_t=\theta_t(e,f)$ be the isosceles triangle with apex angle $\alpha$, base angles $\alpha+\beta$, legs $t\,f\,b$
and base $t\,e\,b$. Its area is $bt^2$ times that of $\tau$. We write $t\in\Theta(e,f)$ if $\theta_t$ is cut into
$bt^2$ copies of $\tau$ (mirror images allowed, not edge-to-edge); in Lean this is
\texttt{ThetaTiled e f b t} (\lean{ThetaUnitStep.lean:38}) with $b=f^2-e^2$: a \texttt{CongruentDissection (b*t*t)} whose
target has sides $(tbf,tbf,tbe)$ and whose model has sides $(ef,b,f^2)$. If $t\in\Theta(e,f)$ then $bt^2$ is a tile
count. For $g\ge1$ let $\alpha_g$ be the isosceles triangle with base angles $\alpha$, legs $gbf^2$ and base
$gb(2f^2-e^2)$ (\texttt{AlphaTiled}).

\paragraph{Prior art (not claimed).} For every tile of the family, $\Theta(e,f)$ contains all large $t$: Paliy's
two-annulus theorem (annuli around the apex adding $e$ and $f$ to the scale from a threshold on). Beeson
(arXiv 1206.2229 v5, Thms.~23, 24, Cor.~4, Fig.~16) gives the $48$-tile seed $\theta_4(1,2)$ and a family of
theta targets. The thresholds below are explicit in $(e,f)$ and the seed is constructed, which is the part we
add; for $e=1$ the threshold is exact.

\subsection{The notch seed family}\label{sub:band:notch}
For coprime $1\le e<f$ put
\[
 n_0=\Big\lceil\frac{e}{f-e}\Big\rceil,\quad j_0=\Big\lceil\frac{n_0e}{b}\Big\rceil,\quad
 M_0=\max\Big(\Big\lceil\frac{fn_0}{e}\Big\rceil,\ n_0+\Big\lceil\frac{en_0}{f}\Big\rceil\Big),\quad
 t_0(e,f)=j_0f+M_0 .
\]

\begin{theorem}[Notch seed family; \sF]\label{thm:notch}
For all coprime $1\le e<f$ and all $t\ge t_0(e,f)$, $t\in\Theta(e,f)$. Hence $(f^2-e^2)t^2$ is a tile count for
every such $t$.
\end{theorem}
\noindent\emph{In words.} Every sufficiently large isosceles triangle $\theta_t(e,f)$ is tiled, and $t_0$ says how large.
\emph{Example.} For $(e,f)=(2,3)$: $n_0=2$, $j_0=1$, $M_0=4$, so $t_0=7$, and $\theta_7(2,3)$ is tiled by $5\cdot49=245$ copies of $\tau(2,3)$.
(For $t=6$ the question is open.)\\
\noindent Lean: \texttt{Erdos634.Draft.ThetaNotch.thetaTiled\_all}
(\lean{ThetaNotch4Rho.lean:156}), with \texttt{t0f} $=t_0$ (\lean{ThetaNotch3Gen.lean:36}); the general case uses
\texttt{thetaTiled\_general} (\lean{ThetaNotch3Gen.lean:99}, hypothesis $e M_0\not\equiv0\bmod f$, proved for
all coprime pairs in \texttt{rho0f\_ne\_zero}) and $e=1$ uses \texttt{thetaTiled\_e1}. The proof is a seed (the
notch frame: a block $(fn)\tau$ at the apex, leg strips, and a notch region filled by one parallelogram of
the lattice of $\tau$) followed by a unit step $\theta_t\to\theta_{t+1}$ (Lean \texttt{ThetaUnit*}); the proofs are in the
earlier draft \rec{notch\_family\_2026-11-14.tex}. $t_0$ is not claimed sharp for $e\ge2$ ($\theta_6(2,3)$ is open
although $t_0(2,3)=7$).

\begin{theorem}[The case $e=1$ is exact; \sF]\label{thm:e1}
For $f\ge2$ and $t\ge1$: $\theta_t(1,f)$ is tiled by copies of $(f,f^2-1,f^2)$ if and only if $t\ge2f$.
\end{theorem}
\noindent\emph{In words.} For $e=1$ the threshold is sharp: $\theta_t(1,f)$ is tiled exactly from $t=2f$ on.
\emph{Example.} $f=2$: the tile is $(2,3,4)$, $\theta_4(1,2)$ has $3\cdot16=48$ tiles (Beeson's $48$-tile seed), and
$\theta_t(1,2)$ is not tiled for $t\le3$.\\
\noindent Lean: \texttt{ThetaNotch.thetaTiled\_e1\_iff} (\lean{ThetaNotchE1.lean:24}), $\texttt{ThetaTiled 1 f (f*f-1) t}\iff 2f\le t$.
The direction ``tiled implies $t\ge2f$'' is \texttt{ThetaKill.theta\_e1}, whose own audit found one inherited gap
(work item \rec{thetakill\_e1\_2026-11-14}, SOUND WITH GAPS); the other direction is Theorem~\ref{thm:notch}.

\begin{proposition}[Alpha hub; \sF]\label{prop:hub}
If $1\le e<f$, $1\le g$ and $t=gf\in\Theta(e,f)$, then $\alpha_g(e,f)$ is tiled: it is the union of $(gb)\tau$ and
$\theta_{gf}$, with $(gb)^2+b(gf)^2=b\,(2f^2-e^2)\,g^2$ tiles.
\end{proposition}
\noindent Lean: \texttt{ThetaNotch.alphaTiled\_of\_thetaTiled} (\lean{ThetaNotch4Glue.lean:178}); instances
\texttt{tileCount\_2576}, \texttt{4025}, \texttt{5796}, \texttt{7008} (\lean{ThetaNotch4Alpha.lean:43ff}). Gluing: Beeson (1206.2229 v5, Cor.~8,
Fig.~25), then Paliy.

\subsection{The band theorem}\label{sub:band:band}
\begin{theorem}[Band theorem; \sF]\label{thm:band}
Let $e,f,n,M,k,s,x,r,m\in\N$ satisfy
\begin{enumerate}
\item[(a)] $1\le e<f$, $1\le n$, $e\le(f-e)\,n$ and $fn\le eM$;
\item[(b)] $0<k<s<M+k$, $e+(f-e)k\le(f-e)s$, and $en+sf\le(M+k)f$;
\item[(c)] $e\mid M+k-s$;
\item[(d)] $e\,k\,(f^2-e^2)=e\,e\,n\,f+x\,(ef)+r\,(f^2-e^2)+m\,f^2$;
\item[(e)] $f\nmid eM$.
\end{enumerate}
Then $M+k\in\Theta(e,f)$; in particular $(f^2-e^2)(M+k)^2$ is a tile count.
\end{theorem}
\noindent\emph{In words.} If the parameters satisfy a short list of inequalities (a frame fits in $\theta_{M+k}$) and
divisibility conditions (the leftover notch region is a union of lattice pieces), then $\theta_{M+k}$ is tiled by an explicit construction.\\
\noindent Lean: \texttt{Erdos634.Draft.ThetaOddE.thetaTiled\_band} (\lean{ThetaOddEBand.lean:33}), with the hypotheses
in the order \texttt{he hef hn ha1 ha2 hk hks hst hd hc hw hY hrho}. Coprimality of $e,f$ is \emph{not} assumed
(a non-coprime pair gives a dilated, still valid, statement); hypothesis (e) is a restriction that the hand
statement did not have (the case $f\mid eM$ is a different case of the notch and is not formalized). Every
quantity is a natural number and all subtractions are of the form $a-b$ with $b\le a$ by the hypotheses.

\paragraph{Proof (sketch; hand proof \sW, formal \sF).} The hand proof is \rec{internal working notes of 2026-11-15}
\S S1--S2 (Theorem~B, proof B.3). The triangle $\theta_{M+k}$ is the union of $k$ leg strips, a triangle $U$, the
block $(fn)\tau$ and a pentagon $R$ (the \emph{notch region}, with corners of angles $\alpha+\beta$, $\gamma$,
$\alpha+\beta$, $\pi-\alpha$, $\pi-\beta$); inequalities (a),(b) say the frame fits; (c),(d) say that the
pentagon is the union of seven convex polygons, each a lattice parallelogram or triangle of $\tau$, whose
sides are integer combinations of $a,b,c$ ($e\mid M+k-s$ for the column widths, (d) for the left wall, which lies on the
base). The Lean proof builds the seven pieces and the volume identity, and obtains the cover from the volume
identity; it does not use the hand text.

\begin{corollary}[Odd $e$; \sF]\label{cor:odde}
For odd $e\ge3$ and $t\ge(e^2+5e+2)/2$: $\theta_t(e,e+1)$ is tiled; hence $(2e+1)t^2$ is a tile count.
\end{corollary}
\noindent\emph{In words.} Tiles $\tau(e,e+1)$ with $e\ge3$ odd have an explicit and simple threshold. \emph{Example.} $e=3$: the
threshold is $(9+15+2)/2=13$, and $\theta_{13}(3,4)$ has $7\cdot169=1183$ tiles.\\
\noindent Lean: \texttt{ThetaOddE.thetaTiled\_oddE} (\lean{ThetaOddEFinal.lean:49}) and \texttt{tileCount\_oddE} (:67). Instances:
$\theta_{13}(3,4)$ (found by machine search on 15 Nov 2026, then built), and the tile counts $756$, $1183$, $2695$,
$77824$, $80275$, $85291$.

\subsection{The $1{,}410$ instances}\label{sub:band:inst}
The modules \texttt{ThetaBandInst1..10}, \texttt{ThetaBandInstH}, \texttt{ThetaBandInstLib}
(\lean{lean/Erdos634/Draft/}) prove \texttt{TileCount N} for $1{,}410$ values $N\le10^6$: $1{,}360$ of the form
$(f^2-e^2)t^2$ by Theorem~\ref{thm:band} or Theorem~\ref{thm:notch} (the parameters are tabulated in each file), and $50$
via Proposition~\ref{prop:hub}. Each instance applies the theorem with literal parameters; the internal audit
(\rec{thetabandinst\_2026-11-15}) read the helper library and sampled instances (SOUND), and a rerun of the census
reproduces exactly this set of $1{,}410$ values.

\subsection{What is not claimed}\label{sub:band:limits}
The band theorem does not say that $\theta_t(e,f)$ is \emph{not} tiled outside its hypotheses; the only exact
threshold is $e=1$. Theorem~\ref{thm:band} reaches $t<t_0(e,f)$ only for some $(e,f,t)$; the table of $t$ below $t_0$
that no frame reaches is part of the open list of Section~\ref{sec:census}. Nothing here bears on the prime case.
\section{Further tile counts by explicit tilings}\label{sec:tilings}
\emph{Roadmap.} One paragraph: what the $132$ tilings are, how they were obtained and checked, and where the details are.
In addition to the theorems above, $132$ integers $N\le10^6$ were shown to be tile counts by explicit tilings
produced by programs and verified by exact rational checkers (mark \sS: machine-verified, \emph{not} checked by the
Lean kernel): $57$ values (tiles with a $2\pi/3$ angle, shapes $F_1,F_2$) obtained by lifting the layout of a known
tiling, $15$ values of shape $F_3$ obtained by Zhang's nesting construction applied to an $F_2$ tiling, and $60$
values of shapes $F_1$ and isosceles obtained by a strip step and a transfer. All three groups are later than
release v5.0, and each tiling was checked by at least two independent programs, with a negative control. The shapes
$F_1,F_2,F_3$ and the nesting are due to Zhang (arXiv 2512.22696) and Harries; the $143$-tiling used in the lift is
Herdt's skeleton; a structural statement of Paliy covers six of the nine $F_2$ values. The strip step and the
transfer are the project's. The constructions, the list of the $132$ values, the verification statement and the
tables are given in the companion data paper; no person outside the project has examined these tilings, and they
are not part of the Lean-verified results of this paper. Their effect on the census is counted in
Section~\ref{sec:census}.

\section{The census of tile counts up to $10^6$}\label{sec:census}

\emph{Roadmap.} We say what the census counts (\ref{sub:cen:def}), give the step-by-step counts and the reconciliation with
release v5.0 (\ref{sub:cen:counts}), describe the four families of open representations, the \emph{walls} (\ref{sub:cen:walls}),
state three open conjectures (\ref{sub:cen:conj}), and say what the census does not show (\ref{sub:cen:not}).

\begin{remark}[Standing]\label{rem:cen:status}
This section reports the output of programs (mark \sS) built on the theorems and tilings of this paper; it is not a
theorem about each integer. The programs are \rec{rust/census\_v5\_20261116} (reproduces the lists below) and the
regression guard \rec{code/census\_guard.sh}; the decision rule for a value excluded is the Lean statement
\texttt{not\_tileCount\_of\_CNzFixW} (Section~\ref{sub:wgap:setting}) or an earlier Lean leaf theorem; a value
included is a tile count by an explicit tiling or a Lean \texttt{TileCount} instance. The
\emph{representation table} of each value (the list of shapes $T,W,\theta,\alpha,\gamma,\dots$ in which it could be a
tile count) is produced by program, and its completeness was re-enumerated independently by program only
(work item \rec{wgap\_gaps\_2026-11-16}, gap~G3): it is not a Lean statement. All audits are internal.
\end{remark}

\subsection{Definition}\label{sub:cen:def}
The \emph{census universe} is the set of non-squarefree integers $N\le10^6$ ($392{,}074$ values). Each $N$ is, in
the census, \emph{decided} (a tile count, by tiling or Lean instance; or not a tile count, by Lean theorems for each
of its representations) or \emph{undecided}. Squarefree $N$ are treated by other theorems of the project and are not
counted here.

\subsection{Counts}\label{sub:cen:counts}
\begin{center}
\begin{tabular}{lrr}
\hline
list (work item \rec{CENSUS\_*}) & undecided $N\le10^6$ & step\\\hline
2026-11-14d & 60,655 & \\
2026-11-15a: band theorem and hub (A3) & 59,245 & $-1{,}410$\\
2026-11-15b: recovered closures and leaf fixes (corrections) & 57,922 & $-921-402$\\
2026-11-15c: isosceles leaf with the second parity (correction) & 57,770 & $-152$\\
2026-11-15d: $F_3$ nesting tilings & 57,755 & $-15$\\
2026-11-16a: W-rung 2 & 57,678 & $-77$\\
2026-11-16b: 57 lift tilings & 57,621 & $-57$\\
2026-11-16c: W-rung 3 & 57,556 & $-65$\\
2026-11-16d: W-rung 4 & 57,458 & $-98$\\
2026-11-16e: W-gap theorem & 52,107 & $-5{,}351$\\
2026-11-16f: two values by \texttt{CNzFixW} & 52,105 & $-2$\\
2026-11-16g: 60 tilings (companion data paper) & \textbf{52,045} & $-60$\\\hline
\end{tabular}
\end{center}
Hence $60{,}655\to52{,}045$ ($-8{,}610$). Of the $392{,}074$ values of the universe, $340{,}029$ are decided in the
census. The undecided values by size: $27$ with $N\le10^3$, $111$ with $N\le3{,}000$, $462$ with $N\le10^4$,
$5{,}115$ with $N\le10^5$, $52{,}045$ with $N\le10^6$. The twenty-seven values $N\le1000$ still open are
\begin{center}
{\small 236, 280, 284, 332, 351, 423, 444, 524, 531, 552, 594, 644, 668, 716, 732, 747, 816, 840, 855,
876, 880, 908, 912, 924, 927, 944, 999.}
\end{center}
Every $N\le200$ is decided (this is unchanged since v4.0 and v5.0).
\emph{Baseline and reconciliation.} The release v5.0 (1 November 2026) published, on the basis of the sifter (non-squarefree
$N\le10^6$), $62{,}372$ undecided values ($129{,}420$ known tile counts, $200{,}282$ excluded). The list of this section
is a subset of the v5.0 list (by a set comparison of the reconstructed v5.0 list with the list 2026-11-16g: no value of the
current list is outside it), so since v5.0 the undecided count fell by exactly $62{,}372-52{,}045=10{,}327$. By category
(one value counted once, at the step that decided it): W-gap theorem $5{,}351$; W rungs $2$--$4$ $77+65+98=240$;
theta and alpha notch family $2{,}694$; band theorem and hub $1{,}410$; recovered closures $319$; isosceles parity $152$;
explicit tilings $161$ ($24+5+15+57+60$; machine-verified, not Lean; the $132=15+57+60$ of Section~\ref{sec:tilings} and $29$ further tilings of
12--14 November 2026, outside the scope of this paper); so
$10{,}166$ values by Lean theorems and hand families and $161$ by tilings. (The chain of the table above starts from the
list of 14 November 2026, which is an earlier state of the sifter and had $1{,}056$ values that v5.0 had already decided;
a difference measured from that start is therefore not a difference since v5.0.) The reconstruction of the v5.0 list is by
set arithmetic on the files of the v5.0 run, not by a rerun of the sifter.

Of the $10{,}327$ values, $471=319+152$ are \emph{corrections} ($319$ closures made earlier and lost by a rebuild of the
program, recovered and now protected by a guard, and $152$ values excluded by a strengthened parity condition on the isosceles
leaf), not new mathematics. The new theorems and families account for the rest.

\subsection{What remains: the four walls}\label{sub:cen:walls}
Every undecided value has at least one open \emph{representation}. The open representations fall into four families,
the \emph{walls}. The sizes below are from the work item tables \rec{SUBATOM\_*\_2026-11-16}; the bases differ (they were
tabulated before the last steps of the table above) and are stated for each line.
\begin{center}
\begin{tabular}{p{4.3cm}p{2.2cm}p{4.2cm}p{2.6cm}}
\hline
wall & open representations & distinct $N$ & base list\\\hline
$T_M(e,f)$ (apex $3\alpha$, base-$\beta$ side) & $26{,}053$ & $20{,}478$ ($19{,}866$ with no other open leaf) & 16f ($52{,}105$)\\
$W_M(e,f)$ (corner $2\alpha$), residual: bottom rung & $2{,}529$ & $2{,}487$ ($2{,}113$ sole) & 15d ($57{,}755$), before W-gap\\
$\theta_t(e,f)$ and the alpha hub & $24{,}089+201$ triples & $17{,}185$ ($16{,}476$ open) & 15d\\
$\gamma=2\alpha$ ($G$) & $35{,}099$ & $26{,}599$ ($10{,}997$ sole) & 15d\\\hline
\end{tabular}
\end{center}
(The $W$ row counts only the bottom-rung representations, since the W-gap representations are now excluded; $W$ before
the theorem had $9{,}239$ open representations, $8{,}785$ values.) The smallest open instance of $T$ is
$236=\tau(4,5)$ at $M=2$; the smallest open $\gamma=2\alpha$ value on its list is $816$ (pair $(7,10)$, $g=4$); the
smallest open $\theta$ cells begin at $N=1836$. On the base list of 2026-11-15c ($57{,}770$, not a release), the leaf sets
were: $T$ alone $19{,}209$; $G+\theta$ $14{,}141$; $G$ $10{,}997$; $W$ $7{,}704$; $\theta$ $2{,}087$; $T+W$ $924$; and
smaller classes, with $2{,}102$ values blocked by a leaf of another type ($E$, $F_1$, Iso, $R_8$).

\subsection{Open statements}\label{sub:cen:conj}
Each of the following would, if proved, decide the stated number of values (counted by the program on the lists named;
\emph{conditional on the representation tables}). None is proved; we do not assert that any of them holds.

\begin{conjecture}[T-threshold; open, \sS-supported only]\label{conj:T}
No representation $T_M(e,f)$ that is open in the program's table is tiled. By the program this would decide $19{,}866$
values with no other open leaf. Evidence on file: engine runs that ended without a tiling on some small instances (observations, not exclusions) and Theorem~\ref{thm:tapex}.
\end{conjecture}
\begin{conjecture}[Theta lower bound; open]\label{conj:theta}
$\theta_t(e,f)$ is not tiled for $t<t_B(e,f)$, the bound tabulated in work item \rec{SUBATOM\_THETA\_2026-11-16}, for the
open cells; this would decide the $\theta$-only values ($2{,}087$) and, together with the next, the $G+\theta$ values.
\end{conjecture}
\begin{conjecture}[Gamma $=2\alpha$; open]\label{conj:G}
Targets $T(g)$ with base angles $\alpha$ of the tiles with $\gamma=2\alpha$ are not tiled for $g\le2v$ (all
representations $(u,v)$, $\gcd(u,v)=1$, $u<v<2u$), beyond the margin $g\le g_D=\lfloor(u-1)/(v-u)\rfloor+1$ that Lean
excludes (Theorem D$'$, \texttt{ThmD.thmD\_all\_g}).
\end{conjecture}
\noindent The bottom rung of $W$ has no conjecture attached: its two corner words are Theorem~\ref{thm:wbottom} and
nothing else is known.

\subsection{What the census does not say}\label{sub:cen:not}
It does not say that the open values are not tile counts; it does not classify the tile counts; the density of
the decided values is a property of the program's tables. It does not touch the prime case, which is settled
by other theorems.

\section*{Credits and acknowledgements for the results of this release}\label{sec:credits_v6}

The results of Sections~\ref{sec:wgap}--\ref{sec:census} rest on the work of others, and we say where.
\begin{itemize}
\item \textbf{M.~Beeson.} The theta targets, the $48$-tile seed $\theta_4(1,2)$ and the hub gluing (arXiv 1206.2229 v5, Thms.~23, 24,
Cor.~4, Cor.~8, Figs.~16, 25); the isosceles-tile program (arXiv 1206.1974, Thm.~11.7, Fig.~24 for the collar at
$(2,3)$); Lemma~52 of 1206.2229 (the finding that the Fig.~5 configuration of Laczkovich is prior art, not a gap).
\item \textbf{M.~Beeson and F.~Ro\c{s}u.} The boundary-run method and Theorem~37 (arXiv 2607.23453), of which
Theorem~\ref{thm:tapex} is the scale-$M$ form at the $3\alpha$ apex, and the corner-run theorem used in
Theorem~\ref{thm:wrun}; the lemmas on the sides of the tile family on which our induction (Lemmas AD, U3, RL) is modelled.
\item \textbf{D.~Paliy.} The first proof of the prime case (2026-09-29), the two-annulus theorem for the isosceles
targets (the large-$t$ statement that Theorem~\ref{thm:notch} makes explicit), the alpha hub and the annulus ideas used for the
band theorem and the collar, and a structural statement covering six of the nine $F_2$ values of Section~\ref{sec:tilings}.
\item \textbf{Y.~X.~Zhang.} The families of tilings with a $2\pi/3$ angle and Props.~9--11 (arXiv 2512.22696): the shapes
$F_1,F_2,F_3$, the nesting of Prop.~11 used for the $F_3$ values, and Props.~10(2), 11(2) (the forty values of the
module \texttt{ZhangInst40}); with Beeson, the rationality results (2604.01314).
\item \textbf{Herdt.} The skeletons: the $143$-tiling (tile $(3,5,7)$), the $1125$ tiling and the layouts of Figs.~23, 24 on which
the lift tool and the collar chain are modelled.
\item \textbf{J.~P.~Harries.} The row constructions $I$--$V$ and the hand tails for the $F_1$ and $F_2$ shapes that
the values of Section~\ref{sec:tilings} are compared with.
\item \textbf{M.~Laczkovich.} The classification of tilings of a triangle (Discrete Math.~140, 1995; Discrete Comput.
Geom.~48, 2012), Lemma~6.2 and Theorem~2.4 of the first and the junction lemmas on which Lemma~M and the thin-$W$
base lemma depend.
\end{itemize}
(The collar construction at $(2,3)$, the instances of Zhang's propositions and the thin-$W$ base lemma were published in
release v5.0 of 1 November 2026 and are not new here.) What is the project's own in this release: the notch seed family and the explicit threshold, the band theorem, the
proofs and Lean statements of the W-gap theorem, its rungs and the structural theorems, the strip step and the
transfer, the census programs and the verification of the tilings. Every Lean statement named in this release is
in the formalization repository with a docstring crediting the sources.

\paragraph{Disclosure.} This work was carried out with substantial assistance from an artificial-intelligence system
(Anthropic's Claude), under the author's direction and review. It has not been refereed, and no reader outside the
project has checked any statement of this release.
\end{document}
